% Source pages for the photographed-page corpus.
%
% Three pages of invented prose: page 1 is the single-page subject, pages 2
% and 3 are the two halves of a spread. Everything here is written for this
% corpus and for nothing else.
%
% Two deliberate departures from ordinary book typesetting, both so that the
% ground-truth text is exactly the text on the page:
%
%   * hyphenation is off, so no word is ever broken across a line;
%   * ligatures are off, so "fi" is two characters on the page and two in the
%     ground truth.
%
% Build: xelatex -output-directory <dir> page.tex

\documentclass[11pt,twoside]{book}

\usepackage[T1]{fontenc}
\usepackage[utf8]{inputenc}
\usepackage[
  paperwidth=5.5in, paperheight=8.5in,
  inner=0.85in, outer=0.7in, top=0.8in, bottom=0.8in,
  headsep=14pt, footskip=26pt
]{geometry}
\usepackage{tgtermes}
\usepackage[verbose=silent]{microtype}
\DisableLigatures[f,q]{encoding = T1}
\usepackage{fancyhdr}

% No hyphenation anywhere: the ground truth is the text, unbroken.
\hyphenpenalty=10000
\exhyphenpenalty=10000
\tolerance=9999
\emergencystretch=3em

\pagestyle{fancy}
\fancyhf{}
\fancyhead[LE]{\small\scshape The Harbour Lights of Calder Bay}
\fancyhead[RO]{\small\scshape Keeping the Lamp}
\fancyfoot[C]{\small\thepage}
\renewcommand{\headrulewidth}{0pt}

\setlength{\parindent}{1.2em}
\setcounter{page}{41}

\begin{document}

%% ---------------------------------------------------------------- page 41
\section*{3\quad Keeping the Lamp}

The keeper who came to Calder Bay in the spring of 1874 found the lamp room
in good order and the stair in none at all. His first report to the board
runs to four pages, of which three concern the stair. It is a document worth
reading, because it shows a man discovering that the work he had been hired
to do was the smaller half of the work in front of him.

A light is not a lamp. The lamp is the least of it: a wick, a reservoir, a
chimney of clear glass. What makes a light is the discipline around the lamp,
and that discipline was written down in a thin book the board issued to every
station on the coast. The book gave the hour of lighting for each week of the
year, the height of the flame in eighths of an inch, and the order in which
the lenses were to be wiped. It did not give reasons. Keepers were expected
to supply those themselves, and most of them did.

Three duties fell outside the book and were understood to be absolute:

\begin{itemize}
  \setlength{\itemsep}{0.15em}
  \item the light burns from sunset to sunrise, in every weather, without
        exception and without shortening;
  \item the log is written the same night, in ink, and never made up from
        memory in the morning;
  \item the store of oil is counted on the first day of each month and the
        count is entered whether or not it agrees with the last one.
\end{itemize}

The third of these caused more trouble than the other two together. Oil
disappeared. It disappeared slowly, in quantities too small to accuse anyone
of, and the board's answer was to demand a count that would make the loss
visible over a season rather than a night. Calder Bay lost eleven gallons in
the year to March 1875, which the keeper recorded without comment, and four
in the year after, which he recorded with a good deal of comment indeed.

\newpage

%% ---------------------------------------------------------------- page 42
What the board never wrote down, and what every keeper learned in the first
autumn, was the difference between a night that is merely dark and a night
that is dark and moving. On a still night the light does its own work. The
keeper trims the wick at ten, again at two, and spends the hours between
them reading or writing or doing nothing at all. On a night of weather he
does not sit down. Salt spray dries on the outer glass in a film that no
amount of rubbing will shift once it has set, so the glass is wiped every
hour whether it needs it or not, and the wiping is done from the inside
where a man can stand.

The Calder Bay tower is eighty-two feet from the rock to the gallery, and
the stair turns eleven times. A keeper who wiped the glass hourly through a
December gale climbed something near two thousand feet in a night, most of
it after midnight, and came off watch at seven to sleep through the short
grey daylight. The board considered this ordinary. So, on the evidence of the
logs, did the keepers.

There is a page in the Calder Bay log for the eighteenth of January 1877 that
has been reproduced more often than any other. The entry is four words long
and reads: \emph{Light burning. Sea over gallery.} The gallery is seventy
feet above high water.

\newpage

%% ---------------------------------------------------------------- page 43
The relief came by boat, and the boat came when it could. In the calmer half
of the year the crossing was an hour and the relief was punctual to the day.
Between November and March it was neither. A keeper due to be relieved on the
first of the month might see the boat on the ninth, or the fourteenth, or not
at all until the weather broke. The station held stores for ninety days for
exactly this reason, and twice in its first decade it needed most of them.

Being late is not the same as being forgotten, and the board took some care
over the difference. Every station on this coast had a signal --- a shape
hoisted on the gallery rail, visible from the headland in clear weather ---
which meant that all was well and no boat was needed. The absence of the
signal meant nothing in itself, because the weather that kept the boat away
also kept the headland out of sight. What the board watched for was the
light. A station that showed its light was a station in working order,
whatever else might be true of it, and a station that failed to show one was
a station in trouble, and the boat went out.

In fifty-one years Calder Bay failed once, for two hours on the night of the
fourth of March 1889, and the reason was entered in the log the same night
in a hand that is noticeably steadier than the sentence it wrote.

\end{document}
